\section{ZK-Rollup Bottleneck Research}
\label{sec:zk_rollup_bottleneck_research}

With the industry shift towards rollups, research has increasingly focused on the specific bottlenecks within the L2 pipeline.

\subsection{Sequencer Policies and Batching}
The sequencer is the central orchestrator of a rollup. Motepalli et al. \cite{motepalli2023sequencers} discuss the implications of centralized sequencers and various transaction ordering policies (e.g., First-Come-First-Served vs. Priority-Fee). The literature highlights that the batch formation strategy directly dictates the trade-off between user latency (waiting for a batch to seal) and L1 cost amortization (waiting for more transactions to share the L1 fee).
\subsection{Prover Overhead}
Generating cryptographic proofs is notoriously resource-intensive. Habib et al. \cite{habib2025bottlenecks} and Chaliasos et al. \cite{chaliasos2024benchmarking, chaliasos2024analyzing} analyze the significant computational overhead introduced by the proving stage. While production systems utilize massive hardware acceleration (GPUs/FPGAs) and recursive proof aggregation, the proving time remains a major contributor to the delay in achieving L1 hard finality.
\subsection{Data Availability (DA) Dynamics}
The cost of posting data to L1 has historically dominated rollup transaction fees. Saif et al. \cite{saif2024dataavailability} and Lavaur et al. \cite{lavaur2023modular} analyze the impact of EIP-4844 blobs. Their work demonstrates that transitioning from standard \texttt{calldata} to blobs significantly reduces DA costs, though it introduces new dynamics regarding blob pricing and temporary data availability windows.
Production ZK systems (e.g., Polygon zkEVM, Scroll, zkSync) use very different hardware setups and prover designs. For instance, Polygon zkEVM demands high-capacity setups (e.g., 96-core CPU, 768 GB RAM) \citep{chaliasos2024analyzing}. This diversity makes identical-hardware comparisons difficult, justifying the use of a research-grade prototype with parameterized proof delays.

\subsection{Summary of Existing Literature}
\label{sec:literature_summary}

Table \ref{tab:literature_comparison} summarizes the key literature reviewed, highlighting the focus areas, methodologies, key findings, limitations, and relevance to this project.

\begin{table}[H]
\centering
\caption{Comparison of Key Literature on Blockchain Scalability and ZK-Rollups}
\label{tab:literature_comparison}
\resizebox{\textwidth}{!}{
\begin{tabular}{|p{3cm}|p{2.5cm}|p{3.5cm}|p{3.5cm}|p{3cm}|p{3cm}|}
\hline
\textbf{Paper / Authors} & \textbf{Focus} & \textbf{Methodology} & \textbf{Key Findings} & \textbf{Limitations} & \textbf{Relevance to this Project} \\ \hline
Chaliasos et al. \citet{chaliasos2024analyzing} & ZK-Rollup Benchmarking & Emprical analysis of public ZK-rollup datasets. & Highlighted L1 DA costs and L2 proving times as major bottlenecks. & Relies on diverse production systems, making apples-to-apples comparisons difficult. & Provides the foundation for the cost models and evaluation metrics used in our controlled framework. \\ \hline
Saif et al. \citet{saif2024dataavailability} & Data Availability in L2s & Survey of DA mechanisms across various rollups. & Identified the shift from calldata to EIP-4844 blobs as a crucial cost-saving measure. & Does not provide empirical execution metrics under varying load. & Motivates our empirical evaluation of Calldata vs. Blob vs. Offchain DA modes. \\ \hline
EIP-4844 \citep{eip4844} & Ethereum Layer-1 Updates & Specification for Proto-Danksharding. & Introduces blob-carrying transactions to drastically lower DA costs. & Theoretical specification; practical L2 integration requires custom sequencer logic. & Informs the implementation of the Blob DA mode and BlobPacking policy. \\ \hline
Sanka et al. \citep{sanka2021systematic} & General Scalability & Systematic literature review of blockchain scaling. & Classifies scaling solutions into L1 and L2, emphasizing the trilemma. & Broad survey; lacks deep technical analysis of ZK proving bottlenecks. & Establishes the broad background and motivation for L2 research. \\ \hline
\end{tabular}
}
\end{table}
